\documentclass[11pt,a4paper]{article}
\usepackage[margin=1in]{geometry}
\usepackage{amsmath}
\usepackage{booktabs}
\usepackage{hyperref}
\usepackage{parskip}

\title{\textbf{Exposé: Multivariate Time Series Forecasting}\\
       \large Deep Learning Bonus Project (SS26) --- TU Darmstadt}
\author{Group ID: TODO \\ Members: TODO}
\date{June 12, 2026}

\begin{document}
\maketitle
\thispagestyle{empty}

\section*{Benchmark Task}

The benchmark dataset contains 96 correlated hourly time series representing
operational load indices across different units.
Each series comes with 28 features including temporal encodings, known-future
covariates (\textit{demand\_forecast}, \textit{queue\_pressure\_forecast},
\textit{staffing\_forecast}, and six others), risk signals, and capacity
metadata.
The task is to forecast 336 hourly steps (14 days) per series, rolled forward
in 24-step blocks.
Evaluation uses MAE, MSE, RMSE, MAPE, sMAPE, and WAPE; lower is better for
all metrics.

\section*{Proposed Architecture: TiDE}

We propose \textbf{TiDE} (Time-series Dense Encoder)~\cite{das2023tide} as our
primary model.
TiDE is an MLP-based encoder-decoder designed specifically for long-horizon
multivariate forecasting.
It encodes a lookback window of past observations together with known-future
covariates through a dense encoder, passes the result through a temporal
decoder that produces one output vector per forecast step, and adds a
residual connection from the input directly to the output.
The full model is implementable in PyTorch using only \texttt{torch.nn} layers,
with no additional dependencies beyond \texttt{torch}, \texttt{pandas}, and
\texttt{numpy}.

\textbf{Why TiDE fits this task.}
The benchmark provides 9 known-future covariate columns that describe
conditions over the full forecast horizon.
TiDE was designed to exploit exactly this signal, which Transformer-based
alternatives such as iTransformer~\cite{liu2024itransformer} handle less
directly.
TiDE is also 5--10$\times$ faster than Transformer baselines~\cite{das2023tide},
which is important given 96 series and a 336-step horizon.

\textbf{Planned improvements and variants.}
We will apply \textbf{RevIN}~\cite{kim2022revin} (Reversible Instance
Normalisation) to handle distribution shift and heterogeneous scales across the
96 series.
We will integrate all 9 known-future covariates as encoder inputs after
imputing missing values via forward-fill.
The lookback window is fixed at \textbf{168 hours} (one full weekly cycle).
Ablations will isolate each component: TiDE without covariates, TiDE with
covariates, and TiDE + RevIN.

\section*{Baselines}

We will compare against the following:
\begin{itemize}
  \setlength\itemsep{0pt}
  \item Naive last-value (minimum required baseline, provided)
  \item Seasonal mean by hour-of-day and day-of-week (provided)
  \item Lag-24 and lag-168 repeat (provided)
  \item Simple LSTM (our neural baseline)
  \item iTransformer~\cite{liu2024itransformer} (ablation comparison)
\end{itemize}

\section*{Additional Dataset: BATADAL}

For the cross-domain experiment we use the \textbf{BATADAL} dataset
(Battle of the Attack Detection Algorithms)~\cite{taormina2018batadal},
a publicly available collection of hourly SCADA sensor readings from a
water distribution network.
Variables include tank water levels (meters) and pipe flow rates (LPS) across
multiple network nodes, recorded continuously over approximately one year.
We filter out labelled attack periods and retain only normal operation rows,
then forecast \textbf{both tank levels and flow rates simultaneously}
(multi-target) using the same 168-hour lookback window.

\textbf{Why BATADAL is suitable.}
Water distribution systems exhibit strong daily and weekly operational cycles,
making them well-suited for deep learning forecasting.
The dataset is rarely used in standard forecasting benchmarks, providing a
genuine novelty contribution.
Thematically, water system operational load is directly analogous to the
benchmark's operational load index, making cross-dataset comparison
well-motivated.
Crucially, BATADAL contains \emph{no} known-future covariates, allowing us to
evaluate whether TiDE generalises from covariate-rich (benchmark) to
covariate-free (BATADAL) settings.
We will report MAE and RMSE on a held-out validation split (last 20\% of
training data).

\section*{Key Risks and Open Questions}

\begin{itemize}
  \setlength\itemsep{0pt}
  \item \textbf{Rolling inference error accumulation:} the model predicts 24
        steps at a time and must roll forward 14 times to cover 336 steps.
        We will evaluate whether re-feeding predictions as history degrades
        performance.
  \item \textbf{Missing covariate values:} several forecast columns contain
        gaps; forward-fill imputation may introduce bias for long gaps.
  \item \textbf{BATADAL preprocessing:} attack periods must be cleanly
        filtered; residual corrupted readings near attack boundaries may affect
        training.
  \item \textbf{Multi-target on BATADAL:} forecasting tank levels and flow
        rates jointly increases output dimensionality; we may fall back to
        single-target if performance suffers.
\end{itemize}

\section*{Main Sources}

\begingroup
\renewcommand{\section}[2]{}
\bibliographystyle{plain}
\bibliography{references}
\endgroup

\end{document}
